\documentclass[11pt, reqno]{amsart}

\usepackage[spanish]{babel}
\usepackage[LGR, T1]{fontenc}
\usepackage[utf8]{inputenc}

../LaTeX/general.tex
% \input{../graphics.tex}

\makeatletter
\def\emailaddrname{\textit{Correo electrónico}}
\def\subtitle#1{\gdef\@subtitle{#1}}
\def\@subtitle{}

% Metadata
\def\logo#1{\gdef\@logo{#1}}
\def\@logo{}
\def\institution#1{\gdef\@institution{#1}}
\def\@institution{}
\def\department#1{\gdef\@department{#1}}
\def\@department{}
\def\professor#1{\gdef\@professor{#1}}
\def\@professor{}
\def\course#1{\gdef\@course{#1}}
\def\@course{}
\def\coursecode#1{\gdef\@coursecode{#1}}
\def\@coursecode{}

\renewcommand{\maketitle}{
\begin{center}
	\small
	\renewcommand{\arraystretch}{1.2}
	\begin{tabular}{cp{.37\textwidth}p{0.44\textwidth}}
		% \hline
		\multirow{5}{*}{\includegraphics[height=2.0cm]{\@logo}}
	  & \multicolumn{2}{c}{ \makecell{{\bfseries \@institution} \\ \@department} } \\
	  % & \multicolumn{2}{c|}{{\bfseries\@institution} \\ \@department} \\
	  \cline{2-3}
	  & \textbf{Profesor:} \@professor & \textbf{Ayudante:} \authors \\
	  % \cline{2-3}
	  & \textbf{Curso:} \@course & \textbf{Sigla:} \@coursecode \\
	  % \cline{2-3}
	  & \multicolumn{2}{l}{ \textbf{Fecha:} \@date } \\
	  % \hline
	\end{tabular}
	\\[\baselineskip]
	% {}
	% \vspace{2\baselineskip}
	{\bfseries\Large\@title}
	\ifx\@subtitle\@empty\else
		\\[1ex]
		\large\mdseries\@subtitle
	\fi
\end{center}
}
\makeatother

\usepackage{multirow, makecell}

\usepackage[
	reversemp,
	letterpaper,
	% marginpar=2cm,
	% marginsep=1pt,
	margin=2.3cm
]{geometry}
\usepackage{fontawesome}
% \makeatletter
% \@reversemargintrue
% \makeatother

% Símbolos al margen, necesitan doble compilación
\newcommand{\hard}{\marginnote{\faFire}}
\newcommand{\hhard}{\marginnote{\faFire\faFire}}

% Dependencias para los teoremas
\usepackage{xifthen}
\def\@thmdep{}
\newcommand{\thmdep}[1]{
	\ifthenelse{\isempty{#1}}
	{\def\@thmdep{}}
	{\def\@thmdep{ (#1)}}
}
\newcommand{\thmstyle}{\color{thm}\sffamily\bfseries}

% ===== Estilos de Teoremas ==========
\newtheoremstyle{axiomstyle}
	{0.3cm}
	{0.3cm}
	{\normalfont}
	{0.5cm}
	{\bfseries\scshape}
	{:}
	{4pt}
	{\thmname{#1}\thmnote{ #3}\thmnumber{ (#2)}}
\newtheoremstyle{styleC}
	{0.5cm}
	{0.5cm}
	{\normalfont}
	{0.5cm}
	{\bfseries}
	{:}
	{4pt}
	{\thmname{#1\textrm{\@thmdep}}\thmnumber{ #2}\thmnote{ (#3)}}

% ====== Teoremas (sin borde) ===========
\theoremstyle{axiomstyle}
\newtheorem*{axiom}{Axioma}

% ====== Teoremas (sin borde) ==================
\theoremstyle{styleC}
\newtheorem{thm}{Teorema}[section]
\newtheorem{mydef}[thm]{Definición}
\newtheorem{prop}[thm]{Proposición}
\newtheorem{cor}[thm]{Corolario}
\newtheorem{lem}[thm]{Lema}
\newtheorem{con}[thm]{Conjetura}
\newtheorem*{sol}{Solución}
\newtheorem*{obs}{Observación}
\newtheorem*{ex}{Ejemplo}

\newtcbox{bluebox}[1][]{enhanced jigsaw, 
  sharp corners,
  frame hidden,
  nobeforeafter,
  listing only,
  #1} % comando para crear cajas de colores

\expandafter\let\expandafter\oldproof\csname\string\proof\endcsname
\let\oldendproof\endproof
\renewenvironment{proof}[1][\proofname]{%
  \oldproof[\scshape Demostración:]%
}{\oldendproof} % comando para redefinir la caja de la demostración
\newenvironment{hint}[1][\proofname]{%
  \oldproof[\scshape Pista:]%
}{\oldendproof} % comando para redefinir la caja de la demostración

% colores utilizados
\definecolor{numchap}{RGB}{249,133,29}
\definecolor{chap}{RGB}{6,129,204}
\definecolor{sec}{RGB}{204,0,0}
\definecolor{thm}{RGB}{106,176,240}
\definecolor{thmB}{RGB}{32,31,31}
\definecolor{part}{RGB}{212,66,66}

% ====== Diseño de los titulares ===============
\usepackage[explicit]{titlesec} % para personalizar el documento, la opción <<explicit>> hace que el texto de los titulares sea un objeto interactuable

\titleformat{\subsection}[runin]
	{\bfseries}
	{\textrm{\S}\thesubsection}
	{1ex}
	{#1.}

\setlist[enumerate,1]{label=\arabic*., ref=\arabic*} % Enumerate standards

\usepackage{tikz}
\usetikzlibrary{babel,cd}
% \DeclareMathOperator\Gr{Gr}

\DeclareMathOperator\pr{pr}

\title{Espacios simplemente conexos}
\date{\DTMdate{2026-10-01}}

\author{José Cuevas Barrientos}
\email{josecuevasbtos@uc.cl}
\urladdr{https://josecuevas.xyz/teach/2026-2-ayud/}

\logo{../puc_negro.png}
\institution{Pontificia Universidad Católica de Chile}
\department{Facultad de Matemáticas}
\course{Topología Algebraica}
\coursecode{MAT2850}
\professor{Mauricio Bustamante}

\begin{document}

\maketitle

\section{Contractibilidad}
\begin{enumerate}
	\item Sea $X$ un espacio arcoconexo.
		Pruebe que las siguientes son equivalentes:
		\begin{enumerate}
			\item $X$ es simplemente conexo.
			\item Dados cualesquiera caminos $\alpha, \beta \colon [0, 1] \to X$ que comparten extremos (i.e., $\alpha(0) =
				\beta(0)$ y $\alpha(1) = \beta(1)$), hay una homotopía (relativa a $\{0, 1 \}$) entre ellos.
			\item Toda función continua $f \colon \SS^1 \to X$ admite una extensión (continua) $\D^2 \to X$ al disco cerrado $\D^2 \subseteq \R^2$.
		\end{enumerate}
		Con ello pruebe que si todo lazo $f \colon \SS^1 \to X$ es homotópico a un lazo constante, pero posiblemente mediante una homotopía no
		relativa o sin punto base.
		Entonces $X$ es simplemente conexo.

	\item\label{ex:easy_van_Kampen}
		Sea $X$ un espacio topológico y sean $U_1, U_2 \subseteq X$ dos abiertos simplemente conexos tales que $X = U_1 \cup U_2$.
		\begin{enumerate}
			\item Pruebe que si su intersección $U_1 \cap U_2 \ne \emptyset$ es arcoconexa, entonces $X$ es simplemente conexo.
			\item\lookst
				Pruebe que si su intersección $U_1 \cap U_2$ tiene exactamente dos componentes arcoconexas, entonces $\pi_1(X,*) \cong \Z$.
		\end{enumerate}

	\item Suponga que $Y$ es contráctil a un punto $y \in Y$.
		Pruebe que para todo espacio topológico $X$ las aplicaciones
		\[
			i \colon X \longrightarrow X \times Y, \qquad x \longmapsto (x, y)
		\]
		y $\pr_1 \colon X\times Y \to X$ son equivalencias homotópicas.

	\item Sea $X$ un espacio arcoconexo.
		\begin{enumerate}
			\item Sea
				\[
					C X := X \times [0,1] / {\sim}
				\]
				su cono, donde la relación de equivalencia $\sim$ identifica $(x,1) \sim (y,1)$.
				Pruebe que $CX$ es contráctil.

			\item\label{ex:suspension_simply_conn}
				Sea 
				\[
					\Sigma X := X \times [-1,1] / {\sim}
				\]
				su suspensión, donde la relación de equivalencia $\sim$ identifica $(x,1) \sim (y,1)$ y $(x,-1) \sim (y,-1)$.
				Pruebe que $\Sigma X$ es simplemente conexo.
		\end{enumerate}
\end{enumerate}

% \section{Ejercicios sugeridos}
% Empleando el cálculo de que $\SS^2$ es simplemente conexo, uno puede probar el siguiente clásico
% \begin{thm}[Borsuk--Ulam]
% 	Para toda función continua $f \colon \SS^n \to \R$ existe un punto $\vec x \in \SS^n \subseteq \R^{n+1}$ tal que $f(\vec x) =
% 	f(-\vec x)$.
% \end{thm}
% \begin{enumerate}[resume]
% 	\item Sean $A_1, A_2, A_3 \subseteq \R^3$ conjuntos compactos (por ende, Lebesgue-medibles de medida finita).
% 		Pruebe que existe un plano $P \subseteq \R^3$ (que no necesariamente pasa por el origen) tal que corta a cada $A_j$ en dos
% 		partes de igual medida.

% 	\item ¿El teorema de Borsuk--Ulam aplica para el toro $\SS^1 \times \SS^1$?
% 		% Dicho de otro modo, toda función continua en $\S$
% \end{enumerate}

\section*{Comentarios}
El ejercicio~\ref{ex:easy_van_Kampen} es un caso sencillo del teorema de Seifert--van Kampen cuyo eslógan es más o menos el siguiente:
\begin{displayquote}
	Si cubrimos un espacio topológico $X$ por abiertos $U_j$ y conocemos el tipo de homotopía de cada $U_j$ y de las intersecciones
	sucesivas; entonces podemos calcular el tipo de homotopía de $X$.
\end{displayquote}
En el ejercicio~\ref{ex:suspension_simply_conn} uno puede preguntarse si $\Sigma X$ es contráctil.
La respuesta es que no siempre, pero no quise incluírlo como ejercicio ya que desconozco si en clases han probado algún ejemplo de un
espacio simplemente conexo que no sea contráctil.

\printbibliography

\end{document}
